% !Mode:: "TeX:UTF-8"
\chapter{绪论}

\section{研究背景与意义}
无人机自主导航通常需要在复杂、动态且部分可观测的环境中完成避障与目标到达任务。在遮挡、光照变化、电磁干扰或传感器老化等条件下，单一传感器容易发生感知退化，导致感知不稳定并进一步影响下游决策。多模态融合（LiDAR、RGB、IMU）能够提供互补信息，但若融合策略无法感知各模态质量变化，则可能在关键时刻放大噪声模态的影响，从而降低系统鲁棒性。

近年来，强化学习（Reinforcement Learning, RL）在端到端控制中表现出较强的策略学习能力，尤其适用于连续动作控制与复杂决策任务。然而，强化学习方法对观测分布漂移和输入噪声较为敏感：当传感器质量发生变化时，策略可能产生性能骤降，甚至引发灾难性碰撞。为此，本文研究一种可解释、可适配的可靠性感知融合机制，并将其与强化学习训练流程紧密集成，旨在提升复杂场景下导航策略的鲁棒性与可用性。

\section{问题定义与研究内容}
\subsection{问题定义}
本文关注的核心问题是：在三模态观测输入（LiDAR点云、RGB图像、IMU惯性数据）存在噪声、缺失或质量波动的条件下，如何设计一种轻量化的可靠性估计与特征融合机制，使融合特征在时变扰动下保持稳定，并能与强化学习算法（SAC）实现端到端协同训练。

\subsection{研究内容}
围绕上述问题，本文的研究内容可归纳为以下四个方面：
\begin{enumerate}
    \item 构建面向三模态输入的质量度量与可靠性估计机制，并给出形式化定义与性质分析；
    \item 设计可端到端训练的特征级融合模块，包括可靠性预测、动态权重分配与特征归一化；
    \item 将融合模块封装为Stable-Baselines3兼容特征提取器，完成与SAC训练流程的工程集成与验证；
    \item 设计实验指标与实验协议，给出消融实验与对比实验计划，并对当前可复现结果进行如实记录。
\end{enumerate}

\section{主要工作与创新点}
本项目（Idea1）在代码仓库中完成实现，核心工作包括：
\begin{enumerate}
    \item \textbf{显式可靠性估计}：针对每种模态构建质量度量模块，输出可解释的可靠性相关特征与评分。
    \item \textbf{可靠性感知自适应融合}：结合可靠性预测网络、基于注意力的动态权重分配以及特征归一化，实现端到端特征级融合。
    \item \textbf{强化学习端到端集成}：将融合模块封装为Stable-Baselines3兼容特征提取器，支持多模态观测字典输入与SAC训练。
\end{enumerate}

\section{技术路线（示意）}
本文的技术路线可概括为"多模态观测 $\rightarrow$ 可靠性估计 $\rightarrow$ 动态加权融合 $\rightarrow$ RL策略学习"。

\textbf{TODO/待补充}：可在后续补充技术路线图（例如系统流程图或模块依赖图），并将仓库中的模块关系（见\texttt{networks/}与\texttt{envs/}）可视化。

\section{论文结构}
本文其余章节组织如下：第2章介绍相关工作与技术背景；第3章给出系统总体设计与方法概述；第4章详细阐述核心模块、关键算法与理论推导；第5章给出实验设计、评价指标与当前可复现实验状态（含TODO占位）；第6章说明工程实现与复现方式；第7章总结全文并展望后续工作。
